\section{Strong-Allee dynamics and analytic basin certificates}
\label{sec:model}

Writing the physical vector state as
\[
  X(t)=(x(t),y(t)),
\]
we consider
\begin{subequations}
\label{eq:model}
\begin{align}
  \CapD x
  &=
  x(1-x)(x-\theta)-axy,
  \label{eq:model-x}\\
  \CapD y
  &=
  y(bx-m),
  \label{eq:model-y}
\end{align}
\end{subequations}
with
\[
  0<\alpha<1,
  \qquad
  a,b,m>0,
  \qquad
  0<\theta<1.
\]
The underlying ecological vector field has been studied previously at integer order
\cite{YeEtAl2019StrongAlleePredatorPrey}; here it serves as a transparent positive system in which continuation-state basin geometry can be resolved rigorously.

If
\[
  \theta<\frac{m}{b}<1,
\]
the positive coexistence equilibrium is
\begin{equation}
  \Eeq=(x^*,y^*),
  \qquad
  x^*=\frac{m}{b},
  \qquad
  y^*=\frac{(1-x^*)(x^*-\theta)}{a}.
  \label{eq:coexistence-equilibrium}
\end{equation}
The Jacobian at \(\Eeq\) is
\begin{equation}
  J=
  \begin{pmatrix}
    x^*(1+\theta-2x^*) & -ax^*\\
    by^* & 0
  \end{pmatrix}.
  \label{eq:coexistence-jacobian}
\end{equation}

\subsection{Cold-start extinction}

\begin{theorem}[Cold-start extinction strip]
\label{thm:extinction-strip}
Assume
\[
  \theta<\frac{m}{b}
\]
and define
\begin{equation}
  \Ext
  =
  \{(x,y):0<x<\theta,\ y\ge0\}.
  \label{eq:extinction-strip}
\end{equation}
Then every standard point initial condition in \(\Ext\) generates a global nonnegative solution satisfying
\[
  X(t)\to(0,0).
\]
Equivalently,
\[
  \iota(\Ext)\subset\Basin(\iota(0,0)).
\]
\end{theorem}

\begin{proof}
The vector field is locally Lipschitz and quasi-positive on the coordinate axes. Caputo viability and extremum results
\cite{GirejkoMozyrskaWyrwas2011Viability,AlRefai2012ExtremePoints}
therefore preserve the nonnegative cone.

Let \(u\) solve
\[
  \CapD u=u(1-u)(u-\theta),
  \qquad
  u(0)=x_0.
\]
Since \(y(t)\ge0\),
\[
  \CapD x
  \le
  x(1-x)(x-\theta).
\]
The scalar comparison principle \cite{Wu2020Comparison} yields
\[
  0\le x(t)\le u(t).
\]
For \(0<x_0<\theta\), the scalar strong-Allee solution remains below \(\theta\) and converges to zero, hence
\[
  x(t)\to0
  \qquad\text{and}\qquad
  x(t)<\theta.
\]

Set
\[
  \delta=m-b\theta>0.
\]
Then
\[
  \CapD y
  =
  y(bx-m)
  \le
  -\delta y.
\]
Comparison with
\[
  \CapD v=-\delta v,
  \qquad
  v(0)=y_0,
\]
gives
\[
  0\le y(t)\le y_0E_\alpha(-\delta t^\alpha)\to0.
\]
The resulting bounds exclude finite-time blow-up by the continuation theorem of Wu and Liu
\cite{WuLiu2020Continuation}. Therefore the solution is global and tends to \((0,0)\).
\end{proof}

\subsection{Inherited-memory survival}

A late physical near-hit to \(\Eeq\) is not a valid survival proof for a Caputo trajectory because restarting at that physical point discards the inherited memory. We therefore retain the entire prehistory explicitly.

Write
\[
  U=X-\Eeq
\]
and
\begin{equation}
  \CapD U=JU+N(U).
  \label{eq:linearized-system}
\end{equation}
Fix a vector norm \(\|\cdot\|\) and its induced matrix norm. Define
\[
  \Psi_J(t)
  =
  t^{\alpha-1}E_{\alpha,\alpha}(Jt^\alpha).
\]
Under the Matignon sector condition, stable Mittag--Leffler estimates imply
\[
  K_J
  :=
  \int_0^\infty\|\Psi_J(s)\|\,ds
  <\infty
\]
\cite{CongDoanTuan2017Perron,CongTuanTrinh2020Asymptotic,Becker2011WeaklySingularResolvents}.

\begin{theorem}[Memory-tail survival criterion]
\label{thm:memory-tail}
Let \(U(t)\) be a solution of \eqref{eq:linearized-system} whose history is known on \([0,T]\). Suppose that, for some \(r>0\),
\[
  \|N(U)\|
  \le
  C_r\|U\|^2
  \qquad
  (\|U\|\le r).
\]
For a finite cut \(T\), define
\begin{equation}
  v_T(t)
  =
  E_\alpha(Jt^\alpha)U_0
  +
  \int_0^T
  \Psi_J(t-s)N(U(s))\,ds,
  \qquad t\ge T,
  \label{eq:memory-response}
\end{equation}
and
\[
  M_T
  =
  \sup_{t\ge T}\|v_T(t)\|.
\]
If
\begin{equation}
  M_T+K_JC_rr^2<r,
  \label{eq:M1}
\end{equation}
then
\[
  U(t)\to0.
\]
\end{theorem}

\begin{proof}
Variation of constants gives the exact split
\begin{equation}
  U(t)
  =
  v_T(t)
  +
  \int_T^t
  \Psi_J(t-s)N(U(s))\,ds.
  \label{eq:exact-memory-split}
\end{equation}
This is a decomposition of the original trajectory, not a restart at time \(T\).

If \(t_e\) were a first exit from the radius-\(r\) ball after \(T\), then
\[
  \|U(t_e)\|
  \le
  M_T+
  \int_T^{t_e}
  \|\Psi_J(t_e-s)\|C_rr^2\,ds
  \le
  M_T+K_JC_rr^2<r,
\]
a contradiction. Hence the trajectory remains in the ball.

Inside the ball,
\[
  \|N(U)\|\le C_rr\|U\|,
\]
and \eqref{eq:M1} implies
\[
  K_JC_rr<1.
\]
The term \(v_T(t)\) tends to zero: the homogeneous term decays by Mittag--Leffler stability, while the finite-history convolution in \eqref{eq:memory-response} tends to zero because \(\Psi_J(t-s)\to0\) for each fixed \(s\in[0,T]\).

Set
\[
  L=\limsup_{t\to\infty}\|U(t)\|.
\]
Splitting the nonlinear convolution in \eqref{eq:exact-memory-split} at a large fixed time and using \(\Psi_J\in L^1(\mathbb{R}_+)\) gives
\[
  L\le K_JC_rr\,L.
\]
Since \(K_JC_rr<1\), one obtains \(L=0\).
\end{proof}
